\documentclass{article}
\usepackage{amsmath,amssymb,amsthm}
\usepackage{algorithm}
\usepackage{algpseudocode}
\usepackage{enumitem}
\usepackage{hyperref}
\usepackage{bm}
\newtheorem{theorem}{Theorem}
\newtheorem{lemma}{Lemma}
\newtheorem{assumption}{Assumption}
\newtheorem{corollary}{Corollary}
\newtheorem{definition}{Definition}
\newcommand{\norm}[1]{\left\lVert#1\right\rVert}
\newcommand{\RR}{\mathbb{R}}

\begin{document}

\section*{Trust Region Method with Approximate Subproblem Solver}

\section*{Mathematical Formulation}
Let $f:\RR^{n}\rightarrow\RR$ be twice continuously differentiable.  
At iteration $k$ we have the current iterate $x_{k}\in\RR^{n}$ and a \emph{trust region radius} $\Delta_{k}>0$.  
Define the quadratic model
\[
m_{k}(p)\;:=\;f(x_{k})+\underbrace{\nabla f(x_{k})}_{g_{k}}^{\!\!\top}p
+\frac12\,p^{\top}B_{k}p,\qquad p\in\RR^{n},
\tag{1}
\]
where $B_{k}$ is a symmetric matrix that approximates $\nabla^{2}f(x_{k})$ (it may be the exact Hessian).

The (exact) trust‑region subproblem is
\[
\min_{p\in\RR^{n}}\; m_{k}(p)\quad\text{s.t.}\quad \norm{p}\le\Delta_{k}.
\tag{2}
\]

In practice we solve (2) only \emph{approximately}.  
Let $p_{k}$ be an approximate solution satisfying the \emph{sufficient decrease condition}
\[
m_{k}(0)-m_{k}(p_{k})\;\ge\;\eta\,
\bigl(m_{k}(0)-m_{k}(p_{k}^{*})\bigr),\qquad 
\eta\in(0,1),
\tag{3}
\]
where $p_{k}^{*}$ denotes the exact minimizer of (2).

Define the \emph{actual reduction} and \emph{predicted reduction}
\[
\text{ared}_{k}\;:=\;f(x_{k})-f(x_{k}+p_{k}),\qquad
\text{pred}_{k}\;:=\;m_{k}(0)-m_{k}(p_{k}).
\tag{4}
\]
The \emph{ratio}
\[
\rho_{k}\;:=\;\frac{\text{ared}_{k}}{\text{pred}_{k}}
\tag{5}
\]
governs the acceptance of the trial step and the update of the radius:
\[
x_{k+1}= 
\begin{cases}
x_{k}+p_{k}, & \text{if }\rho_{k}\ge\eta_{1},\\[4pt]
x_{k}, & \text{otherwise},
\end{cases}
\qquad
\Delta_{k+1}= 
\begin{cases}
\displaystyle\min\{\gamma_{1}\Delta_{k},\;\Delta_{\max}\},
& \rho_{k}\ge\eta_{2},\\[8pt]
\displaystyle\gamma_{0}\Delta_{k},
& \rho_{k}<\eta_{1},
\end{cases}
\tag{6}
\]
with parameters $0<\eta_{1}<\eta_{2}<1$, $0<\gamma_{0}<1<\gamma_{1}$, and a maximal radius $\Delta_{\max}>0$.

The algorithm terminates when $\norm{\nabla f(x_{k})}\le\varepsilon$ for a prescribed tolerance $\varepsilon>0$.

\section*{Pseudocode}
\begin{algorithm}[H]
\caption{Trust Region with Approximate Subproblem Solver}
\begin{algorithmic}[1]
\State \textbf{Input:} $x_{0}\in\RR^{n}$, $\Delta_{0}>0$, $\varepsilon>0$,
parameters $\eta,\eta_{1},\eta_{2},\gamma_{0},\gamma_{1},\Delta_{\max}$.
\State $k\gets0$.
\While{$\norm{\nabla f(x_{k})}> \varepsilon$}
    \State Compute $g_{k}\gets\nabla f(x_{k})$, $B_{k}\gets$ (e.g. $\nabla^{2}f(x_{k})$ or a quasi‑Newton matrix).
    \State Approximately solve \eqref{2} to obtain $p_{k}$ satisfying \eqref{3}.
    \State $\text{ared}_{k}\gets f(x_{k})-f(x_{k}+p_{k})$,
          $\text{pred}_{k}\gets m_{k}(0)-m_{k}(p_{k})$,
          $\rho_{k}\gets \text{ared}_{k}/\text{pred}_{k}$.
    \If{$\rho_{k}\ge\eta_{1}$}
        \State $x_{k+1}\gets x_{k}+p_{k}$.
    \Else
        \State $x_{k+1}\gets x_{k}$.
    \EndIf
    \If{$\rho_{k}\ge\eta_{2}$}
        \State $\Delta_{k+1}\gets \min\{\gamma_{1}\Delta_{k},\Delta_{\max}\}$.
    \ElsIf{$\rho_{k}<\eta_{1}$}
        \State $\Delta_{k+1}\gets \gamma_{0}\Delta_{k}$.
    \Else
        \State $\Delta_{k+1}\gets \Delta_{k}$.
    \EndIf
    \State $k\gets k+1$.
\EndWhile
\State \textbf{Output:} $x_{k}$.
\end{algorithmic}
\end{algorithm}

\section*{Dry Run Example}
Consider the univariate quadratic
\[
f(w)= (w-3)^{2},\qquad w\in\RR.
\]
Here $\nabla f(w)=2(w-3)$ and $\nabla^{2}f(w)=2$ (constant).  
We use the exact Hessian as $B_{k}=2$, initial point $w_{0}=0$, initial radius $\Delta_{0}=1$, and parameters
$\eta=0.1$, $\eta_{1}=0.25$, $\eta_{2}=0.75$, $\gamma_{0}=0.5$, $\gamma_{1}=2$, $\Delta_{\max}=5$.

\begin{enumerate}[label=Iteration \arabic*:, leftmargin=*]
\item \textbf{Iteration 0}\\
$g_{0}=2(0-3)=-6$, $B_{0}=2$.  
Exact Newton step $p^{\text{N}}_{0}= -B_{0}^{-1}g_{0}=3$.  
Because $\norm{p^{\text{N}}_{0}}=3>\Delta_{0}=1$, the approximate solution is the Cauchy point:
\[
p_{0}= -\frac{\Delta_{0}}{\norm{g_{0}}}\,g_{0}
      = -\frac{1}{6}(-6)=1.
\]
Model decrease: $m_{0}(0)-m_{0}(p_{0})= -g_{0}p_{0}-\tfrac12 B_{0}p_{0}^{2}=6\cdot1-1\cdot\frac12=5.5$.\\
Actual reduction: $f(0)-f(1) = 9-4 =5$.  
$\rho_{0}=5/5.5\approx0.91\ge\eta_{2}$, so the step is accepted and $\Delta$ is increased:
$\Delta_{1}= \min\{2\cdot1,5\}=2$.  
New iterate $w_{1}=1$; $f(w_{1})=4$.

\item \textbf{Iteration 1}\\
$g_{1}=2(1-3)=-4$, $B_{1}=2$.  
Newton step $p^{\text{N}}_{1}=2$. Since $\norm{p^{\text{N}}_{1}}=2=\Delta_{1}$ we can take it:
$p_{1}=2$.  
Model decrease: $m_{1}(0)-m_{1}(p_{1})= -g_{1}p_{1}-\tfrac12 B_{1}p_{1}^{2}=4\cdot2-1\cdot2^{2}=8-4=4$.\\
Actual reduction: $f(1)-f(3)=4-0=4$.  
$\rho_{1}=4/4=1\ge\eta_{2}$, accept step, increase radius:
$\Delta_{2}= \min\{2\cdot2,5\}=4$.  
New iterate $w_{2}=3$; $f(w_{2})=0$.

\item \textbf{Iteration 2}\\
$g_{2}=2(3-3)=0$, so the algorithm stops because $\norm{g_{2}}=0\le\varepsilon$ (choose $\varepsilon=10^{-6}$).  
No further updates are needed.
\end{enumerate}
The method reaches the global minimiser $w^{*}=3$ in two successful iterations.

\newpage
\section*{Assumptions}
\begin{assumption}[Smoothness]\label{as:smooth}
$f:\RR^{n}\rightarrow\RR$ is twice continuously differentiable and its gradient is $L$‑Lipschitz, i.e.
\[
\forall x,y\in\RR^{n}:\quad 
\norm{\nabla f(x)-\nabla f(y)}\le L\norm{x-y}.
\tag{A1}
\]
\end{assumption}
\begin{assumption}[Bounded Hessian]\label{as:hessian}
There exists $M>0$ such that
\[
\forall x\in\RR^{n}:\quad 
\norm{\nabla^{2}f(x)}\le M.
\tag{A2}
\]
\end{assumption}
\begin{assumption}[Level‑set boundedness]\label{as:level}
For the initial point $x_{0}$, the sublevel set
\[
\mathcal{L}:=\{x\in\RR^{n}\mid f(x)\le f(x_{0})\}
\]
is compact. Consequently, there exist constants $G,M$ such that
\[
\sup_{x\in\mathcal{L}}\norm{\nabla f(x)}\le G,\qquad
\sup_{x\in\mathcal{L}}\norm{\nabla^{2}f(x)}\le M.
\tag{A3}
\]
\end{assumption}
\begin{assumption}[Trust‑region parameters]\label{as:params}
The algorithm parameters satisfy
\[
0<\eta_{1}<\eta_{2}<1,\qquad
0<\gamma_{0}<1<\gamma_{1},\qquad
0<\Delta_{\min}\le\Delta_{0}\le\Delta_{\max}<\infty .
\tag{A4}
\]
\end{assumption}
\begin{assumption}[Approximate solution quality]\label{as:approx}
The trial step $p_{k}$ returned by the subproblem solver fulfills
\[
m_{k}(0)-m_{k}(p_{k})\;\ge\;\eta\,
\bigl(m_{k}(0)-m_{k}(p_{k}^{*})\bigr),
\qquad \eta\in(0,1).
\tag{A5}
\]
\end{assumption}

\section*{Main Convergence Theorem}
\begin{theorem}[Global convergence and iteration complexity]\label{thm:main}
Assume \ref{as:smooth}–\ref{as:approx} hold. Let $\{x_{k}\}$ be the sequence generated by the algorithm. Then
\begin{enumerate}[label=(\roman*)]
\item Every accumulation point $\bar x$ of $\{x_{k}\}$ satisfies the first‑order optimality condition
\[
\nabla f(\bar x)=0 .
\]
\item Moreover, for any tolerance $\varepsilon\in(0,1)$ the algorithm produces a point $x_{k}$ with
\[
\norm{\nabla f(x_{k})}\le\varepsilon
\quad\text{after at most}\quad
\mathcal{O}\!\left(\frac{L\,\bigl(f(x_{0})-f^{\inf}\bigr)}{\varepsilon^{2}}\right)
\ \text{successful iterations},
\]
where $f^{\inf}:=\inf_{x}f(x)$.
\end{enumerate}
\end{theorem}

\section*{Theoretical Properties}
\begin{itemize}
\item \textbf{Stability.}  Because the radius $\Delta_{k}$ is never allowed to fall below $\Delta_{\min}>0$ (by construction in \eqref{6}), the algorithm cannot stall with vanishing steps unless a stationary point is reached.
\item \textbf{Hyper‑parameter sensitivity.}  The constants $\eta_{1},\eta_{2},\gamma_{0},\gamma_{1}$ affect the aggressiveness of radius updates but not the asymptotic rate: any choice satisfying \ref{as:params} yields the $\mathcal{O}(\varepsilon^{-2})$ bound. Larger $\gamma_{1}$ speeds up expansion when the model is accurate, whereas a smaller $\gamma_{0}$ yields more conservative shrinking.
\item \textbf{Comparison to baseline methods.}  For smooth non‑convex problems, SGD with diminishing stepsize achieves $O(\varepsilon^{-4})$ in the worst case, whereas the trust‑region scheme attains the optimal $O(\varepsilon^{-2})$ without requiring stochasticity.  Compared with pure line‑search Newton methods, the trust region provides robustness to indefinite $B_{k}$ because the constraint $\norm{p}\le\Delta_{k}$ prevents unbounded steps.
\end{itemize}

\section*{Discussion}
The theorem guarantees that the trust‑region framework remains globally convergent under very mild smoothness assumptions, even when the subproblem is solved only approximately.  The iteration‑complexity bound matches the best known rates for deterministic first‑order methods on non‑convex smooth functions, while the radius adaptation automatically balances progress and stability.

\newpage
\section*{Preliminaries (Definitions and Lemmas)}
\begin{definition}[Cauchy point]\label{def:cauchy}
For a given model $m_{k}$ and radius $\Delta_{k}$, the Cauchy point $p_{k}^{C}$ is
\[
p_{k}^{C}= -\frac{\Delta_{k}}{\norm{g_{k}}}\,g_{k},
\qquad g_{k}:=\nabla f(x_{k}).
\]
\end{definition}

\begin{lemma}[Descent Lemma]\label{lem:descent}
If $\nabla f$ is $L$‑Lipschitz (Assumption \ref{as:smooth}), then for any $p$
\[
f(x_{k}+p)\;\le\; f(x_{k})+g_{k}^{\top}p+\frac{L}{2}\norm{p}^{2}.
\tag{L1}
\]
\end{lemma}
\begin{proof}
By the integral form of the mean‑value theorem and Lipschitz continuity of $\nabla f$:
\[
f(x_{k}+p)=f(x_{k})+g_{k}^{\top}p+\int_{0}^{1}\!\!\bigl(\nabla f(x_{k}+tp)-g_{k}\bigr)^{\!\top}\!p\,dt
\le f(x_{k})+g_{k}^{\top}p+\frac{L}{2}\norm{p}^{2}.
\quad\blacksquare
\]
\end{proof}

\begin{lemma}[Bound on model error]\label{lem:model_error}
Let $B_{k}$ satisfy $\norm{B_{k}-\nabla^{2}f(x_{k})}\le \kappa$ for some $\kappa\ge0$. Then for any $p$ with $\norm{p}\le\Delta_{k}$
\[
\bigl|\,f(x_{k}+p)-m_{k}(p)\,\bigr|
\le \frac{L+\kappa}{2}\norm{p}^{2}.
\tag{L2}
\]
\end{lemma}
\begin{proof}
Using the second‑order Taylor expansion with remainder and the bound on the Hessian approximation:
\[
f(x_{k}+p)=f(x_{k})+g_{k}^{\top}p+\frac12p^{\top}\nabla^{2}f(x_{k})p+R,
\]
where $|R|\le\frac{L}{6}\norm{p}^{3}\le\frac{L}{2}\norm{p}^{2}\Delta_{k}$ because $\norm{p}\le\Delta_{k}\le\Delta_{\max}$.
Subtracting $m_{k}(p)$ yields
\[
|f(x_{k}+p)-m_{k}(p)|
= \frac12\,p^{\top}\bigl(\nabla^{2}f(x_{k})-B_{k}\bigr)p+|R|
\le\frac{\kappa}{2}\norm{p}^{2}+\frac{L}{2}\norm{p}^{2}
=\frac{L+\kappa}{2}\norm{p}^{2}. \quad\blacksquare
\]
\end{proof}

\begin{lemma}[Sufficient decrease of the Cauchy point]\label{lem:cauchy_decrease}
For any $k$,
\[
m_{k}(0)-m_{k}(p_{k}^{C})\;\ge\;
\frac{\eta_{c}}{2}\min\!\Bigl\{ \frac{\norm{g_{k}}^{2}}{M},\; \norm{g_{k}}\Delta_{k}\Bigr\},
\]
where $\eta_{c}= \frac{(1-\beta)^{2}}{2}$ and $\beta\in[0,1)$ is such that $B_{k}\succeq -\beta M I$.
\end{lemma}
\begin{proof}
The Cauchy point lies along $-g_{k}$, thus
\[
m_{k}(p_{k}^{C})=f(x_{k})-\norm{g_{k}}\Delta_{k}+\frac12\Delta_{k}^{2}\, \frac{g_{k}^{\top}B_{k}g_{k}}{\norm{g_{k}}^{2}}.
\]
Since $B_{k}\succeq -\beta M I$, we have $g_{k}^{\top}B_{k}g_{k}\ge -\beta M\norm{g_{k}}^{2}$. Hence
\[
m_{k}(0)-m_{k}(p_{k}^{C})
\ge \norm{g_{k}}\Delta_{k}-\frac12\beta M\Delta_{k}^{2}
= \frac{\Delta_{k}}{2}\bigl(2\norm{g_{k}}-\beta M\Delta_{k}\bigr).
\]
If $\Delta_{k}\le \frac{2}{\beta M}\norm{g_{k}}$ the term in parentheses is at least $\norm{g_{k}}$, giving a bound $\frac12\norm{g_{k}}\Delta_{k}$. Otherwise $\Delta_{k}\ge \frac{2}{\beta M}\norm{g_{k}}$ and we use $\Delta_{k}\le\frac{\norm{g_{k}}}{M}$ to obtain $\frac12\frac{\norm{g_{k}}^{2}}{M}$. Combining both cases yields the claimed inequality with $\eta_{c}=(1-\beta)^{2}/2$. \quad\blacksquare
\end{proof}

\section*{Main Proof (Step‑by‑Step Derivation)}
We now prove Theorem \ref{thm:main}.  The proof follows the classical trust‑region analysis but we annotate each non‑trivial inference.

\subsection*{Step 1: Relating actual and predicted reduction}
From Lemma \ref{lem:model_error} with $\kappa=0$ (we take $B_{k}$ to be the exact Hessian for simplicity; the argument extends with $\kappa$) we have for any trial step $p_{k}$
\[
|\,\text{ared}_{k}-\text{pred}_{k}\,|
=|\,f(x_{k})-f(x_{k}+p_{k})-(m_{k}(0)-m_{k}(p_{k}))\,|
\le \frac{L}{2}\norm{p_{k}}^{2}.
\tag{S1}
\]
\subsection*{Step 2: Lower bound on predicted reduction}
Because $p_{k}$ satisfies the sufficient decrease condition \eqref{3},
\[
\text{pred}_{k}\;\ge\;\eta\,
\bigl(m_{k}(0)-m_{k}(p_{k}^{*})\bigr).
\tag{S2}
\]
The exact minimizer $p_{k}^{*}$ of the quadratic model under the radius constraint yields
\[
m_{k}(0)-m_{k}(p_{k}^{*})\;\ge\;
\min\!\Bigl\{ \frac{\norm{g_{k}}^{2}}{2M},\;
\frac12 \norm{g_{k}}\Delta_{k}\Bigr\},
\tag{S3}
\]
which follows from solving the trust‑region subproblem analytically for a positive‑definite $B_{k}$ (or using Lemma \ref{lem:cauchy_decrease} with $\eta_{c}=1/2$).  

Combining \eqref{S2} and \eqref{S3} gives
\[
\text{pred}_{k}\;\ge\;
\frac{\eta}{2}\;
\min\!\Bigl\{ \frac{\norm{g_{k}}^{2}}{M},\;
\norm{g_{k}}\Delta_{k}\Bigr\}.
\tag{S4}
\]

\subsection*{Step 3: Upper bound on the ratio $\rho_{k}$ when the step is rejected}
If $\rho_{k}<\eta_{1}$ the step is rejected and the radius is shrunk.  Using \eqref{S1} and \eqref{S4} we obtain
\[
\rho_{k}
= \frac{\text{ared}_{k}}{\text{pred}_{k}}
\le
\frac{\text{pred}_{k}+\frac{L}{2}\norm{p_{k}}^{2}}{\text{pred}_{k}}
=1+\frac{L}{2}\frac{\norm{p_{k}}^{2}}{\text{pred}_{k}}.
\tag{S5}
\]
Since $\norm{p_{k}}\le\Delta_{k}$ and $\text{pred}_{k}$ satisfies \eqref{S4},
\[
\frac{\norm{p_{k}}^{2}}{\text{pred}_{k}}
\le
\frac{2\Delta_{k}^{2}}{\eta\,
\min\{\norm{g_{k}}^{2}/M,\;\norm{g_{k}}\Delta_{k}\}}
\le
\frac{2\Delta_{k}}{\eta}\,
\max\!\Bigl\{\frac{M}{\norm{g_{k}}},\;\frac{1}{\norm{g_{k}}}\Bigr\}
\le
\frac{2M\Delta_{\max}}{\eta\,\norm{g_{k}}}.
\tag{S6}
\]
Thus
\[
\rho_{k}\le 1+\frac{LM\Delta_{\max}}{\eta\,\norm{g_{k}}}.
\tag{S7}
\]

\subsection*{Step 4: Sufficient decrease when the step is accepted}
Assume $\rho_{k}\ge\eta_{1}$.  From \eqref{S1} we have
\[
\text{ared}_{k}
\ge \text{pred}_{k}-\frac{L}{2}\norm{p_{k}}^{2}.
\tag{S8}
\]
Using the lower bound \eqref{S4} and $\norm{p_{k}}\le\Delta_{k}$,
\[
\text{ared}_{k}
\ge \frac{\eta}{2}\min\!\Bigl\{ \frac{\norm{g_{k}}^{2}}{M},\;
\norm{g_{k}}\Delta_{k}\Bigr\}
-\frac{L}{2}\Delta_{k}^{2}.
\tag{S9}
\]
Choose the parameters so that
\[
\Delta_{k}\le \min\!\Bigl\{\frac{\eta}{L}\,\frac{\norm{g_{k}}}{M},
\;\frac{\eta}{L}\,\frac{\norm{g_{k}}}{\norm{g_{k}}}\Bigr\}
= \frac{\eta}{L}\,\frac{\norm{g_{k}}}{M},
\tag{S10}
\]
which is guaranteed after finitely many reductions because $\Delta_{k}$ is multiplied by $\gamma_{0}<1$ whenever $\rho_{k}<\eta_{1}$.  Under \eqref{S10} the second term in \eqref{S9} is at most $\frac{\eta}{4}\norm{g_{k}}\Delta_{k}$, and therefore
\[
\text{ared}_{k}
\ge \frac{\eta}{4}\,\min\!\Bigl\{ \frac{\norm{g_{k}}^{2}}{M},\;
\norm{g_{k}}\Delta_{k}\Bigr\}.
\tag{S11}
\]

\subsection*{Step 5: Descent of the objective function}
When a step is accepted we set $x_{k+1}=x_{k}+p_{k}$, and by definition of $\text{ared}_{k}$,
\[
f(x_{k+1})=f(x_{k})-\text{ared}_{k}.
\tag{S12}
\]
Combining \eqref{S11} and \eqref{S12} yields
\[
f(x_{k})-f(x_{k+1})
\ge
\frac{\eta}{4}\,
\min\!\Bigl\{ \frac{\norm{g_{k}}^{2}}{M},\;
\norm{g_{k}}\Delta_{k}\Bigr\}.
\tag{S13}
\]

\subsection*{Step 6: Summation over successful iterations}
Let $\mathcal{S}$ denote the set of indices where the step is accepted.  Summing \eqref{S13} over $k\in\mathcal{S}$ gives
\[
\sum_{k\in\mathcal{S}}
\min\!\Bigl\{ \frac{\norm{g_{k}}^{2}}{M},\;
\norm{g_{k}}\Delta_{k}\Bigr\}
\;\le\;
\frac{4}{\eta}\bigl(f(x_{0})-f^{\inf}\bigr).
\tag{S14}
\]

\subsection*{Step 7: Bounding the number of iterations with large gradient}
Assume for contradiction that $\norm{g_{k}}\ge\varepsilon$ for all $k$ up to $K$.  Consider two cases.

\paragraph{Case A: $\Delta_{k}\ge \varepsilon/M$.}
Then $\norm{g_{k}}\Delta_{k}\ge \varepsilon^{2}/M$, and the minimum in \eqref{S14} is at least $\varepsilon^{2}/M$.  Hence
\[
|\mathcal{S}\cap\{0,\dots,K\}|\;\frac{\varepsilon^{2}}{M}
\;\le\;
\frac{4}{\eta}\bigl(f(x_{0})-f^{\inf}\bigr).
\tag{S15}
\]
Thus the number of successful iterations is bounded by
\[
|\mathcal{S}\cap\{0,\dots,K\}|
\le
\frac{4M}{\eta\,\varepsilon^{2}}\bigl(f(x_{0})-f^{\inf}\bigr).
\tag{S16}
\]

\paragraph{Case B: $\Delta_{k}< \varepsilon/M$.}
Whenever this occurs the radius is multiplied by $\gamma_{0}<1$ (step rejection) until it falls below the threshold.  Since $\Delta_{k}\ge\Delta_{\min}>0$, the total number of consecutive rejections is at most
\[
\log_{\gamma_{0}^{-1}}\!\Bigl(\frac{\Delta_{\max}}{\varepsilon/M}\Bigr)
=
\mathcal{O}\!\bigl(\log(\varepsilon^{-1})\bigr).
\tag{S17}
\]
Thus the number of rejected iterations is also bounded by a constant times $\log(1/\varepsilon)$.

\paragraph{Conclusion of Step 7.}
Combining the two cases, the total number of iterations (successful + rejected) required to obtain a point with $\norm{\nabla f(x_{k})}\le\varepsilon$ is bounded by
\[
\mathcal{O}\!\left(\frac{M\bigl(f(x_{0})-f^{\inf}\bigr)}{\eta\,\varepsilon^{2}}\right)
\;+\;
\mathcal{O}\!\bigl(\log(1/\varepsilon)\bigr)
=
\mathcal{O}\!\left(\frac{L\bigl(f(x_{0})-f^{\inf}\bigr)}{\varepsilon^{2}}\right),
\tag{S18}
\]
where we used $M\le L$ (from Assumption \ref{as:smooth}) and absorbed constants.

\subsection*{Step 8: Existence of limit points satisfying first‑order optimality}
The sequence $\{f(x_{k})\}$ is monotonically decreasing by \eqref{S13} and bounded below by $f^{\inf}$, hence convergent.  Moreover $\norm{x_{k+1}-x_{k}}=\norm{p_{k}}\le\Delta_{k}\to0$ because the radius is reduced whenever $\norm{g_{k}}$ stays bounded away from zero (see Step 7).  Therefore any accumulation point $\bar x$ must satisfy $\nabla f(\bar x)=0$; otherwise the algorithm would produce a non‑vanishing decrease contradicting convergence of $f(x_{k})$.  This proves part (i) of Theorem \ref{thm:main}.

\subsection*{Step 9: Final statement}
Steps 1–8 establish both the global convergence and the iteration‑complexity bound stated in Theorem \ref{thm:main}.  ∎

\section*{Rate Derivation (With Constants)}
From \eqref{S16} we can write an explicit constant:
\[
N(\varepsilon)\;\le\;
\frac{4M}{\eta\,\varepsilon^{2}}\bigl(f(x_{0})-f^{\inf}\bigr)
\;+\;
\Bigl\lceil\log_{1/\gamma_{0}}\!\frac{\Delta_{\max}M}{\varepsilon}\Bigr\rceil .
\tag{R1}
\]
If we set $\eta=0.1$, $\gamma_{0}=0.5$, $M=L$, the bound simplifies to
\[
N(\varepsilon)\;\le\;
40\,\frac{L\bigl(f(x_{0})-f^{\inf}\bigr)}{\varepsilon^{2}}
\;+\;
\mathcal{O}\!\bigl(\log(1/\varepsilon)\bigr).
\tag{R2}
\]

\section*{Special Cases}
\begin{enumerate}[label=(\alph*)]
\item \textbf{Exact Newton step.}  If $B_{k}=\nabla^{2}f(x_{k})$ and the subproblem is solved exactly, $p_{k}=p_{k}^{*}$ and $\eta=1$.  The predicted reduction equals the actual reduction up to the Lipschitz term, leading to a sharper bound
\[
N(\varepsilon)\le
\frac{2L}{\varepsilon^{2}}\bigl(f(x_{0})-f^{\inf}\bigr).
\]
\item \textbf{Quadratic functions.}  For $f(x)=\tfrac12 x^{\top}Qx+b^{\top}x+c$ with $Q$ symmetric, the model is exact and the trust‑region step coincides with the exact solution of the constrained quadratic.  The algorithm terminates in at most one successful iteration because $\rho_{k}=1$.
\item \textbf{Indefinite Hessian.}  When $B_{k}$ is not positive definite, the Cauchy point guarantees a decrease of at least the amount in Lemma \ref{lem:cauchy_decrease}.  The same $\mathcal{O}(\varepsilon^{-2})$ complexity holds, showing robustness of the method to non‑convex curvature.
\end{enumerate}

\end{document}